\documentclass[10pt,a4paper]{amsart}
\usepackage[margin=19mm]{geometry}
\usepackage{amsmath,amssymb,mathtools}
\usepackage[scr=rsfs,cal=euler]{mathalfa}
\usepackage{xfrac}
\usepackage[unicode,hidelinks,bookmarks=false]{hyperref}
\newcommand{\ie}{i.e.\ }
\newcommand{\cf}{cf.\ }
\newcommand{\resp}{respectively\ }
\newcommand{\Subsection}{\subsection}
\providecommand{\nobreakdash}{\nobreak\hbox{-}}
\setlength{\emergencystretch}{2em}
\multlinegap0pt
\usepackage{xcolor}
\usepackage{amssymb}
\usepackage{mathtools}
\usepackage[shortlabels]{enumitem}
\newtheorem{proposition}{Proposition}
\newcommand{\tr}{\operatorname{tr}}
\title{Clifford Realizations of Hyperbolic Cubics}
\author{Tim Netzer}
\date{Source: arXiv:2609.12957v1 (CC BY 4.0)}
\begin{document}
\maketitle

\noindent\emph{Source and licence.} Clifford Realizations of Hyperbolic Cubics. Version arXiv:2609.12957v1 (CC BY 4.0), distributed by arXiv under the Creative Commons Attribution 4.0 International licence, http://creativecommons.org/licenses/by/4.0/, as recorded in the arXiv licence metadata for this exact version and shown on the abstract page https://arxiv.org/abs/2609.12957v1. Full licence notice: \texttt{licenses/arxiv-2609.12957-CC-BY-4.0.txt}.\par\smallskip

\noindent\textit{Editorial scope.} Counted unit: the article's proposition prop:contact-data (source0.tex lines 1380--1411). The article's complete original proof of it is delivered with it (source0.tex lines 1413--1529, one \texttt{proof} environment of 117 lines). No mathematics is written, repaired or re-derived: the statement and the proof are the article's own lines, in the article's own notation, and no retained line is substituted or re-flowed.  No further source-local context is needed: every cross-reference of every retained line resolves inside the two delivered spans.  Nothing was omitted from any retained span, so the package carries no omission record.  No retained line contains a citation, so the wrapper carries no bibliography and no bibliography entry.  No retained line contains a figure environment, an image-inclusion command or drawing code, so no image is delivered and the package declares no asset dependency.  Editorial formatting only, in addition: an AMS \texttt{amsart} preamble that restates the article's own declarations and macros used by the retained lines, namely the environments \texttt{proposition} on separate counters, together with the standard packages those lines need.  The retained environments print in this document as Proposition 1 in order of appearance, whereas the article prints the same items with the numbering of its own full document.  The complete native source of the article as distributed in the arXiv e-print for this exact version is bundled unchanged as \texttt{sources/210165/source0.tex}.
\begin{proposition}
\label{prop:contact-data}
Let $q\in H_4$ satisfy $q(e_1)=1$, and suppose that $h_q$ admits a
quaternionic Hermitian determinantal representation
\[
 h_q(t,x)=\det_M(tI_3+A(x)),\qquad
 A\colon\mathbb R^4\to\operatorname{Herm}_3(\mathbb H)
 \text{ linear}.
\]
Put $U=e_1^\perp$.  Since $q$ attains its maximum on the unit sphere
at $e_1$, its derivative in every direction of $U$ vanishes there.
Thus the term quadratic in $s$ and linear in $y$ is zero, and we may write
\[
 q(s,y)=s^3+3sy^*Dy+r(y),
 \qquad s\in\mathbb R,\quad y\in U,
\]
where $D=D^*$ and $r$ is cubic.  Let $T$ be the canonical symmetric
pencil of $r$, and put
\[
 P=\frac{I+D}{2},\qquad Q=\frac{I-D}{2}.
\]
Then $Q\geqslant I/4$, and there are a finite-dimensional real Hilbert
space $F$, a unit vector $\xi\in F$, a linear map $a\colon U\to F$, and
a linear map $J\colon U\to\operatorname{Hom}(U,F)$ such that
\[
 \begin{gathered}
 a^*\xi=0,\qquad J(y)^*\xi=Dy,\qquad a^*a=P,\\
 a^*J(y)+J(y)^*a=T(y),\qquad
 J(y)^*J(y)=\|y\|^2Q.
 \end{gathered}
\]
\end{proposition}

\begin{proof}
Comparing coefficients in $\det_M(tI_3+A(x))=h_q(t,x)$ gives
\[
 \tr A(x)=0,\qquad
 \frac12\tr(A(x)^2)=3\|x\|^2,
 \qquad \det_M A(x)=2q(x).
\]
Since $h_q(t,e_1)=(t-1)^2(t+2)$, a quaternionic unitary conjugation
lets us assume $A(e_1)=\operatorname{diag}(2,-1,-1)$.
For $y\in U$, let $\alpha(y)$ and $B(y)$ be the diagonal blocks of
$A(0,y)$.  Taking the trace of $A(0,y)$ and comparing the coefficients
of $s$ in $\frac12\tr(A(s,y)^2)=3(s^2+\|y\|^2)$ give, respectively,
$$\alpha(y)+\tr B(y)=2\alpha(y)-\tr B(y)=0.$$
Thus $\alpha=0$, $\tr B=0$, and, for a linear map $z\colon U\to\mathbb H^2$,
\[
 A(s,y)=
 \begin{pmatrix}
  2s&z(y)^*\\
  z(y)&-sI_2+B(y)
 \end{pmatrix},
 \qquad y\in U.
\]
Since $B(y)\in\operatorname{Herm}_2(\mathbb H)$ is traceless, there is a
self-adjoint $G\geqslant0$ on $U$ such that
$B(y)^2=(y^*Gy)I_2$.  Taking the trace of the square and expanding the
Moore determinant by the lower-right Schur complement give the following
identities.  The inverse simplifies because $B(y)^2$ is scalar, and
continuity covers singular blocks:
\[
 \|z(y)\|^2+y^*Gy=3\|y\|^2,
 \qquad
 q(s,y)=s^3+\frac{s}{2}
 \bigl(\|z(y)\|^2-2y^*Gy\bigr)
 +\frac12z(y)^*B(y)z(y).
\]
Comparing this with the prescribed form of $q$ and extracting mixed terms gives
\begin{equation}\label{eq:defect-identities}
 \begin{aligned}
 G&=I-2D,\\
 B(y)B(u)+B(u)B(y)&=2(y^*Gu)I_2,\\
 \operatorname{Re}\bigl(z(u)^*z(v)\bigr)
   &=u^*(3I-G)v,\\
 r(y)&=\frac12z(y)^*B(y)z(y).
 \end{aligned}
\end{equation}
The third identity gives $G\leqslant3I$.  Thus
$Q=(I+G)/4\geqslant I/4$ and $P=(3I-G)/4$.
Define the real-bilinear maps $L,\mathcal A\colon U\times U\to\mathbb H^2$ by
\[
 \begin{aligned}
 L(y,u)&=\frac16B(y)z(u)+\frac13B(u)z(y),\\
 \mathcal A(y,u)&=B(y)z(u)-B(u)z(y).
 \end{aligned}
\]
In orthonormal coordinates on $U$, put
\[
 w(y,u)=(y_1u_2-y_2u_1,\ y_1u_3-y_3u_1,\ y_2u_3-y_3u_2)^T.
\]
Coordinate expansion, using the alternation of $\mathcal A$, shows that
the right-hand side below is quadratic in these three minors appearing in $w$.  Define
the real symmetric matrix $K$ by
\begin{equation}\label{eq:K-coordinate}
 w(y,u)^*Kw(y,u)=\frac1{18}\|\mathcal A(y,u)\|^2
 +(y^*Dy)(u^*Du)-(y^*Du)^2.
\end{equation}
Every unit vector in $\mathbb R^3$ is $w(u_1,u_2)$ for an orthonormal
pair.  A plane rotation preserves this vector and makes $u_1^*Gu_2=0$.
Put $g_i=u_i^*Gu_i\in[0,3]$ and
\[
 X=B(u_1)z(u_2),\qquad Y=B(u_2)z(u_1),\qquad
 \alpha=\sqrt{(3-g_1)(3-g_2)},\quad \beta=2\sqrt{g_1g_2}.
\]
By \eqref{eq:defect-identities},
$\|X\|\|Y\|=\alpha\beta/2$.  Expanding
\eqref{eq:K-coordinate} and applying Cauchy--Schwarz therefore gives
\[
 \begin{aligned}
 36\,w(u_1,u_2)^*Kw(u_1,u_2)
 &=\alpha^2+\beta^2-4\operatorname{Re}(X^*Y)\\
 &\geqslant\alpha^2+\beta^2-2\alpha\beta
 =(\alpha-\beta)^2\geqslant0.
 \end{aligned}
\]
Thus $K\geqslant0$.  Take the real Hilbert space
$F=\mathbb R\oplus\mathbb H^2\oplus\mathbb R^3$, where the quaternionic
summand carries the real inner product
$\langle p,q\rangle=\operatorname{Re}(p^*q)$.
The adjoints of $a$ and $J(y)$ below are taken with respect to these real
Hilbert space structures.  Set $\xi=(1,0,0)$ and define
\[
 \begin{aligned}
 a(u)&=(0,\tfrac12z(u),0),\\
 J(y)u&=(y^*Du,L(y,u),K^{1/2}w(y,u)).
 \end{aligned}
\]
The identities $a^*\xi=0$ and $J(y)^*\xi=Dy$ are immediate, while
$a^*a=P$ follows from
$\operatorname{Re}(z(u)^*z(v))=4u^*Pv$ in \eqref{eq:defect-identities}.
Extracting mixed terms from $r$ gives
\[
 \frac12\operatorname{Re}\bigl(z(u)^*L(y,v)
 +L(y,u)^*z(v)\bigr)
 =u^*T(y)v.
\]
Hence $a^*J(y)+J(y)^*a=T(y)$.  Finally, cancellation of the mixed terms
and \eqref{eq:defect-identities} give
\[
 \begin{aligned}
 \|L(y,u)\|^2+\frac1{18}\|\mathcal A(y,u)\|^2
 &=\frac1{12}\|B(y)z(u)\|^2+\frac16\|B(u)z(y)\|^2\\
 &=\|y\|^2u^*Qu-(y^*Dy)(u^*Du).
 \end{aligned}
\]
Together with \eqref{eq:K-coordinate}, this gives
$\|J(y)u\|^2=\|y\|^2u^*Qu$ for every $u$, hence
$J(y)^*J(y)=\|y\|^2Q$.
\end{proof}

\end{document}
